\section{Our Method}
\label{sec:OurMethod}

\subsection{Preliminaries}

\textbf{Privileged information setup.} Following the LUPI paradigm~\cite{vapnik2009new,lopez2016unifying}, we treat ED as \emph{privileged information}: it is computed via DFT during training but is not required at inference, eliminating the cost of quantum chemical computation at deployment.

\textbf{Notation and Problem Formulation.} The molecular geometry and the corresponding ground-truth labels with $k$ prediction tasks are $\{\mathcal{G}_i=(\mathcal{V}_i, \mathcal{E}_i)\}_{i=1}^n$ and $\{y_i\}_{i=1}^n\in\mathbb{R}^k$ respectively, where $\mathcal{V}_i\in\mathbb{R}^{n^v_i\times(3+d_i^v)}$ ($n^v_i$, 3, $d_i^v$ are the number of atoms, the coordinates, and the feature dimensions) and $\mathcal{E}_i\in\mathbb{R}^{n^e_i\times d_i^e}$ ($n^e_i$ is the number of edges and $d_i^e$ is the feature dimensions of bonds). The corresponding multi-view ED image and structural image are $\mathcal{U}\in\mathbb{R}^{N_v\times 4\times H\times W}$ and $\mathcal{S}\in\mathbb{R}^{N_s\times 3\times H\times W}$ respectively, where $N_v$ and $N_s$ denote the number of views for ED images ($N_v\!=\!6$) and structural images ($N_s\!=\!4$), and $H$, $W$ are the image height and width. 3 and 4 represent RGB and RGB-D channels, respectively.

The framework comprises three stages: (1) pre-train a masked autoencoder (MAE)~\cite{he2022masked} with ED encoder $f_{EDE}$ and decoder $f_{EDD}$ to learn representations $\mathcal{F}^{\mathcal{U}}\in\mathbb{R}^{d^U}$ from ED images $\mathcal{U}$; (2) jointly pre-train an ED-aware teacher $f_{S}$, ED predictor $f_{EDP}$, and reliability estimator $\phi_{eval}$ to map structural images $\mathcal{S}$ to structural features $\mathcal{F}^{\mathcal{S}}\in\mathbb{R}^{d^{S}}$, then to ED-related features $\mathcal{F}^{\mathcal{S}\rightarrow\mathcal{U}}\in\mathbb{R}^{d^U}$, and to predict the teacher's per-sample reconstruction error; (3) distill a geometry student $f_{G}$ using the ED-aware teacher, ED predictor, reliability estimator, and mapper $f_M$ with selective distillation.

\subsection{Overview of the Method}

We propose \textbf{EDG} (\textbf{E}lectron \textbf{D}ensity-enhanced \textbf{G}eometry learning) and its reliability-aware extension \textbf{EDG++}. As shown in Fig.~\ref{Fig2_Method}, the framework has four modules: (a) 2 million molecular conformations are processed with DFT to produce multi-view ED images with RGB-D channels (Section~\ref{sec:Generation_ED_Image}); (b) ImageED learns ED-related representations from these images via two self-supervised tasks (Section~\ref{sec:ED_representation_learning}); (c) an ED-aware teacher is trained to predict ED features from structural images by minimizing the discrepancy with ImageED features (Section~\ref{sec:Pretraining_of_ED-aware_Teacher}); (d) the teacher distills ED knowledge into a geometry student through an ED predictor and mapper (Section~\ref{sec:ED-enhanced_Molecular_Geometry_Learning}). EDG++ extends the base framework with a theoretical analysis of when selective distillation reduces gradient noise (Section~\ref{sec:Theoretical_Analysis}), a reliability estimator (Section~\ref{sec:Reliability_Estimator}), and an adaptive thresholding mechanism (Section~\ref{sec:Adaptive_Thresholding}). Algorithm~\ref{alg:edg_process} in Appendix~\ref{sec:appendix_algorithm} summarizes the full three-stage pipeline.

\begin{figure*}[!htb]
\centering
\includegraphics[width=\textwidth]{imgs_edgpp/our_method/EDG++框架图_en.png}
\caption{Overview of the proposed EDG/EDG++ framework. \textbf{(a)} Multi-view ED images with RGB-D channels are generated based on 2 million molecular conformers and DFT data, which contains the structural and ED information of the molecule. \textbf{(b)} ImageED with masked autoencoder architecture and 2 pretext tasks ($\mathcal{L}_{VR}$ and $\mathcal{L}_{MR}$) is pretrained to extract ED-related features from multi-view ED images in (a). \textbf{(c)} The ED-aware teacher takes the structural images as input and learns to transform them into ED features $\mathcal{F}^{\mathcal{S}\rightarrow\mathcal{U}}$ using ED predictor and ED encoder in (b). The reliability estimator $\phi_{eval}$ is trained to predict the reconstruction error. \textbf{(d)} In downstream tasks, the ED-aware teacher, ED predictor, and reliability estimator from (c) are frozen to enhance the geometry student. EDG uses naive distillation while EDG++ uses reliability-aware selective distillation with adaptive thresholding. Electron density is no longer required at this stage, as the frozen ED-aware teacher already encodes ED-related information from structural images alone.}
\label{Fig2_Method}
\end{figure*}

\subsection{Generation of RGB-D Electron Density Images}
\label{sec:Generation_ED_Image}

We adopt the multi-view RGB-D rendering protocol of EDG~\cite{xiang2025edg}, summarized here for self-containment. As shown in Fig.~\ref{Fig2_Method}(a), 2 million conformations from PCQM4Mv2~\cite{hu2021ogb} are processed using DFT calculations~\cite{turney2012psi4} to obtain ED and electrostatic potential (ESP) grids, then rendered via PyMol~\cite{delano2002pymol} into 6-view RGB-D images $\mathcal{U}\in\mathbb{R}^{6\times 4\times 224\times 224}$: RGB encodes the ESP via a red-white-blue colormap, and depth encodes spatial layout. Details of DFT settings, PyMol commands, and viewpoint configurations are in Appendix~\ref{sec:appendix_pymol}.

This image format produces a fixed-size output tensor regardless of source ED grid resolution and, in our prior work~\cite{xiang2025edg}, improved energy prediction by 38.4\%, GPU memory efficiency by 42.1\%, and training efficiency by 4.8\% over the best point-cloud / voxel alternatives on a 10K-molecule benchmark; per-task RMSE values are reproduced in Table~\ref{tab:results_ED_images} (Appendix~\ref{sec:appendix_ed_repr}). Qualitative ImageED reconstruction examples are provided in Appendix~\ref{sec:appendix_imageed_vis}.

\subsection{ED Representation Learning with ImageED}
\label{sec:ED_representation_learning}

We train ImageED, the ED feature extractor used by the teacher in Section~\ref{sec:Pretraining_of_ED-aware_Teacher}, on the 2M ED images from Section~\ref{sec:Generation_ED_Image}. ImageED follows the ViT-Base/16~\cite{dosovitskiy2020image} masked-autoencoder backbone of EDG~\cite{xiang2025edg}: each multi-view ED image is patchified (patch size 16, $14\times 14$ grid per view, 25\% mask ratio), encoded by an ED encoder $f_{EDE}$ on visible tokens, and decoded by $f_{EDD}$ on the union of encoded and learnable mask tokens, yielding reconstructed unmasked / masked tokens $\hat{t}^u$ / $\hat{t}^m$.

% [Polish] Original: ... visible reconstruction on unmasked tokens (renamed from EDG's $\mathcal{L}_{MP}$ to $\mathcal{L}_{VR}$ for alignment ...
Two complementary objectives capture local and contextual ED structure: \emph{visible reconstruction} on unmasked tokens (renamed from $\mathcal{L}_{MP}$ in EDG to $\mathcal{L}_{VR}$ for alignment with standard MAE terminology and to avoid notation conflict with the new selective-distillation \emph{mask} introduced in Section~\ref{sec:Adaptive_Thresholding}), and \emph{masked reconstruction} on masked tokens (the standard MAE objective~\cite{he2022masked}, renamed to $\mathcal{L}_{MR}$ for symmetry):
\begin{align}
\label{Formula_L_VR}
\mathcal{L}_{VR}=\frac{1}{B}\sum_{i=1}^B d(t^u_i, \hat{t}^u_i),\quad
\mathcal{L}_{MR}=\frac{1}{B}\sum_{i=1}^B d(t^m_i, \hat{t}^m_i)
\end{align}
where $d(\cdot,\cdot)$ denotes Euclidean distance. The overall objective is
\begin{align}
\label{Formula_L_ImageED}
\mathcal{L}_{ImageED}=\lambda_{VR}\mathcal{L}_{VR}+\lambda_{MR}\mathcal{L}_{MR},\quad \lambda_{VR}=\lambda_{MR}=1,
\end{align}
following~\cite{xiang2025edg}. After pre-training with 2 million molecules, $f_{EDE}$ is frozen and used to extract ED features for downstream stages.

\subsection{Pretraining of ED-aware Teacher}
\label{sec:Pretraining_of_ED-aware_Teacher}

Since computing ED for every downstream task is prohibitive, we learn the mapping from conformations directly to ED features, bypassing explicit ED computation. The generation path is: conformation $\mathcal{S}\xrightarrow{\text{DFT}}$ ED data $\mathcal{U}\xrightarrow{\text{ImageED}}$ ED features $\mathcal{F}^{\mathcal{U}}$. We adopt a Markov assumption on this chain (\ie, $\mathcal{F}^{\mathcal{U}}$ depends on $\mathcal{S}$ only through $\mathcal{U}$), which justifies learning $p(\mathcal{F}^{\mathcal{U}}|\mathcal{S})$ without ED data at inference (see Appendix~\ref{sec:appendix_markov}). Under the Born-Oppenheimer framework for ground-state systems with fixed composition and charge, the Hohenberg-Kohn theorem guarantees that the electron density is uniquely determined by the nuclear geometry, supporting this assumption. In practice, the lossy multi-view projection from 3D density fields to 2D images introduces residual dependence that the Markov assumption does not capture; the consistent improvements across four architectures on two datasets suggest that this approximation is adequate.

We instantiate this mapping with an ED-aware teacher $f_{\mathcal{S}}$ and ED predictor $f_{EDP}$. The teacher uses ResNet18~\cite{he2016deep} with view-wise average pooling to extract structural features from 4-view structural images $\mathcal{S}$ rendered via PyMol. The ED predictor $f_{EDP}$ is a 2-layer MLP (Linear$\rightarrow$Softplus$\rightarrow$Linear) that maps structural features to ED features:
\begin{align}
\label{Formula_F_S_U_and_F_S}
\mathcal{F}^{\mathcal{S}\rightarrow\mathcal{U}}=f_{EDP}(\mathcal{F}^{\mathcal{S}}); \mathcal{F}^{\mathcal{S}}=f_{\mathcal{S}}(s)
\end{align}

We freeze the pre-trained ED encoder and use token-wise average pooling to obtain ED features $\mathcal{F}^{\mathcal{U}}$ as the training target. The alignment loss is:
\begin{align}
\label{Formula_L_align}
\mathcal{L}_{align}=L1(\mathcal{F}^{\mathcal{S}\rightarrow\mathcal{U}}, \mathcal{F}^{\mathcal{U}})
\end{align}
where $L1$ denotes L1 distance. With the ED-aware teacher, the costly ED image can be replaced by a cheaper structural image, eliminating DFT-related costs in downstream tasks.

\subsection{Theoretical Analysis of Selective Distillation}
\label{sec:Theoretical_Analysis}

We analyze why selective distillation can outperform naive distillation when teacher reliability varies across samples. The naive distillation gradient combines a \emph{signal} component (pulling the student toward true ED features) with a \emph{noise} component (pulling toward the teacher's reconstruction error); when teacher errors are heavy-tailed, filtering the highest-error subset reduces noise faster than it reduces the sample count. Proposition~\ref{prop:gradient_noise} formalizes this intuition as a sufficient condition under simplifying assumptions (bounded Jacobians, independence of $\|\varepsilon_i\|$ and $\|J_i\|_F$), yielding an upper-bound result that the empirical improvements (Section~\ref{sec:Experiments}) and tail-error evidence (Section~\ref{sec:error_distribution}) complement. The reliability estimator (Section~\ref{sec:Reliability_Estimator}) and adaptive thresholding mechanism (Section~\ref{sec:Adaptive_Thresholding}) implement this filtering in practice.

\textbf{Noisy Teacher Model.} In cross-modal distillation, the teacher's ED feature prediction for sample $i$ can be decomposed as:
\begin{equation}
\label{eq:teacher_decomp}
    \mathcal{F}^{\mathcal{S}\rightarrow\mathcal{U}}_i = \mathcal{F}^{\mathcal{U}}_i + \varepsilon_i
\end{equation}
% [Polish] Original: where ... incur significant information loss.
where $\mathcal{F}^{\mathcal{U}}_i \in \mathbb{R}^{d^U}$ is the ground-truth ED feature from ImageED and $\varepsilon_i \in \mathbb{R}^{d^U}$ is the reconstruction error of the teacher. In cross-modal settings, $\|\varepsilon_i\|$ varies substantially across samples: some molecular geometries are faithfully captured by the multi-view projection, while others (\eg, molecules with delocalized electronic structure or symmetric charge distributions) incur substantial information loss.

\textbf{Distillation Gradient Decomposition.} Using the squared L2 loss $\|\cdot\|^2$ for analytical tractability (the qualitative conclusions extend to smooth L1 via local quadratic approximation), the naive distillation gradient for a batch $\mathcal{B}$ decomposes as:
\begin{equation}
\label{eq:grad_decomp}
    \nabla_\theta \mathcal{L}^{naive}_{ED} = \underbrace{\frac{2}{|\mathcal{B}|} \!\sum_{i \in \mathcal{B}} \!\left(\mathcal{F}^{\mathcal{G}\!\rightarrow\!\mathcal{U}}_i - \mathcal{F}^{\mathcal{U}}_i\right)^{\!\top} \!\! J_i}_{\displaystyle g_{\text{signal}}} \; \underbrace{- \frac{2}{|\mathcal{B}|} \!\sum_{i \in \mathcal{B}} \varepsilon_i^\top J_i}_{\displaystyle g_{\text{noise}}}
\end{equation}
where $J_i = \nabla_\theta \mathcal{F}^{\mathcal{G}\rightarrow\mathcal{U}}_i$ is the Jacobian of the student's predicted ED features. The noise term $g_{\text{noise}}$ deflects optimization away from the true ED features. In cross-modal distillation, $\varepsilon_i$ is generally not zero-mean---the teacher \emph{systematically} errs on certain molecular conformations---causing consistent directional bias. Our upper-bound analysis (Proposition~\ref{prop:gradient_noise}) holds regardless of the mean structure of $\varepsilon_i$, as it relies on Cauchy--Schwarz rather than cancellation of cross terms.

\begin{proposition}[Gradient Noise Reduction via Selective Distillation]
\label{prop:gradient_noise}
Let $\mathcal{B}$ be a training batch and $\mathcal{S} \subseteq \mathcal{B}$ be a selected subset that excludes samples with the largest teacher errors $\|\varepsilon_i\|$. Define $\bar{r}_{\mathcal{S}} = \frac{1}{|\mathcal{S}|}\sum_{i \in \mathcal{S}} \|\varepsilon_i\|^2$ and $\bar{r}_{\mathcal{B}} = \frac{1}{|\mathcal{B}|}\sum_{i \in \mathcal{B}} \|\varepsilon_i\|^2$ as the mean squared teacher errors. Under the assumptions that (i) $\|\varepsilon_i\|$ and $\|J_i\|_F$ are independent, and (ii) $\|J_i\|_F \leq C$ for all $i$, upper bounds on the expected squared gradient noise satisfy:
\begin{equation}
\label{eq:noise_ratio}
    \frac{\mathbb{E}\left[\|g^{\mathcal{S}}_{\text{noise}}\|^2\right]}{\mathbb{E}\left[\|g^{\mathcal{B}}_{\text{noise}}\|^2\right]} \leq \frac{|\mathcal{B}|}{|\mathcal{S}|} \cdot \frac{\bar{r}_{\mathcal{S}}}{\bar{r}_{\mathcal{B}}}
\end{equation}
Consequently, selective distillation reduces the expected gradient noise whenever:
\begin{equation}
\label{eq:improvement_condition}
    \frac{\bar{r}_{\mathcal{S}}}{\bar{r}_{\mathcal{B}}} < \frac{|\mathcal{S}|}{|\mathcal{B}|}
\end{equation}
\end{proposition}

The proof bounds each term under the independence assumption (Appendix~\ref{sec:appendix_proof}). Proposition~\ref{prop:gradient_noise} assumes oracle access to true teacher errors $\|\varepsilon_i\|$ for subset selection; in practice EDG++ uses a learned reliability estimator (Section~\ref{sec:Reliability_Estimator}) as a proxy, and the achievable bound tightens as the estimator's ranking correlates with $\|\varepsilon_i\|$. Condition~\eqref{eq:improvement_condition} holds whenever discarded samples carry \emph{disproportionately} high noise relative to their count---\ie, when teacher errors are heavy-tailed---giving a sufficient (not tight) condition for noise reduction under the bounded-Jacobian and Cauchy--Schwarz simplifications.

Section~\ref{sec:error_distribution} provides \emph{indirect but supportive} evidence for the heavy-tailed structure. The evidence is indirect because we observe downstream student prediction errors rather than the teacher reconstruction errors $\|\varepsilon_i\|$ in Proposition~\ref{prop:gradient_noise}. The two quantities are causally linked: samples with large $\|\varepsilon_i\|$ receive noisier supervision, producing systematically larger prediction errors. This pattern holds across three architectures and four property types (27--38\% tail error reduction), supporting the causal interpretation. The independence assumption between $\|\varepsilon_i\|$ and $\|J_i\|_F$ is also conservative: molecular structural complexity simultaneously increases both teacher error and student sensitivity, inducing positive dependence that tightens the bound (Appendix~\ref{sec:appendix_remarks}).

\subsection{Reliability Estimator}
\label{sec:Reliability_Estimator}

\noindent\textit{Nomenclature.} Throughout this paper we use \emph{reliability score} (or simply \emph{reliability}) for the scalar $c_i = -\phi_{eval}(\mathcal{F}^{\mathcal{S},(i)})$, with higher values indicating higher estimated reliability of the teacher's per-sample prediction. The term ``confidence'' is reserved for statistical quantities such as bootstrap confidence intervals; ``trust'' and ``fidelity'' appear only in their natural-language sense and do not denote separate quantities.

Following Proposition~\ref{prop:gradient_noise}, we introduce a reliability estimator $\phi_{eval}: \mathbb{R}^{d^S} \rightarrow \mathbb{R}$ as a practical proxy for the per-sample teacher error $\|\varepsilon_i\|$. It is a lightweight 2-layer MLP (Linear(512, 256)$\rightarrow$Softplus$\rightarrow$Linear(256, 1)) jointly trained with the ED-aware teacher in Stage~II. Both its input features and regression target are detached from the teacher's computation graph to prevent gradient interference: without input detachment, the evaluator's gradient could push the teacher toward making errors more \emph{predictable} rather than more \emph{accurate}; without target detachment, $\mathcal{L}_{eval}$ would back-propagate through $\mathcal{F}^{\mathcal{S}\rightarrow\mathcal{U}}$ and alter the alignment loss landscape. The training objective is:
\begin{equation}
    \mathcal{L}_{eval} = \mathbb{E}\left[\left|\phi_{eval}\!\bigl(\operatorname{sg}(\mathcal{F}^{\mathcal{S}})\bigr) - \operatorname{sg}\!\bigl(\|\mathcal{F}^{\mathcal{S}\rightarrow\mathcal{U}} - \mathcal{F}^{\mathcal{U}}\|_2\bigr)\right|\right]
\end{equation}
where $\operatorname{sg}(\cdot)$ denotes stop-gradient (detach).
At inference, we convert the evaluator output to a reliability score $c_i = -\phi_{eval}(\mathcal{F}^{\mathcal{S},(i)})$, where higher $c_i$ indicates higher estimated reliability. The threshold form $\tau = \mu + \kappa\sigma$ (Section~\ref{sec:Adaptive_Thresholding}) is inherently scale-invariant, ensuring that $\kappa$ controls the retention rate in units of standard deviations regardless of the absolute scale of $c_i$. Although introduced as a training-time gate, the same frozen estimator can be queried at inference to flag per-sample OOD inputs (\eg, MD frames drifting from the training manifold), a capability exploited in the Chignolin case study (Sec.~\ref{sec:case_study_chignolin}).

\subsection{Adaptive Thresholding Mechanism}
\label{sec:Adaptive_Thresholding}

We implement selective distillation with an adaptive thresholding mechanism combining global (dataset-level) and local (batch-level) statistics. The form $\tau = \mu + \kappa \cdot \sigma$ is \emph{distribution-agnostic}: since $\mu$ and $\sigma$ come from the same reliability score population, a given $\kappa$ yields comparable filtering strictness regardless of the score distribution's scale or shape. No distributional assumption is required; under approximate normality, the retention rate simplifies to $\Phi(-\kappa)$ (Appendix~\ref{sec:appendix_remarks}).

\textbf{Local Threshold.} For each mini-batch $\mathcal{B}$, we compute:
\begin{equation}
    \tau_{local} = \mu_\mathcal{B} + \kappa_{batch} \cdot \sigma_\mathcal{B}
\end{equation}
where $\mu_\mathcal{B}$ and $\sigma_\mathcal{B}$ are the mean and standard deviation of reliability scores within the batch, and $\kappa_{batch}$ controls the strictness of filtering.

\textbf{Global Threshold.} We pre-compute statistics over the entire training set:
\begin{equation}
    \tau_{global} = \mu_{all} + \kappa_{all} \cdot \sigma_{all}
\end{equation}

\textbf{Hybrid Threshold.} The final threshold combines both:
\begin{equation}
    \tau = \beta \cdot \tau_{local} + (1 - \beta) \cdot \tau_{global}
\end{equation}
where $\beta \in [0, 1]$ balances local adaptivity and global consistency. Setting $\beta = 0$ yields pure global thresholding, $\beta = 1$ yields pure local thresholding, and $\beta = 0.5$ provides a balanced hybrid strategy.

\subsection{ED-enhanced Molecular Geometry Learning}
\label{sec:ED-enhanced_Molecular_Geometry_Learning}

During downstream training, the geometry student $f_{G}$ and the frozen ED-aware teacher extract features $\mathcal{F}^{\mathcal{G}}$ and $\mathcal{F}^{\mathcal{S}}$ from geometry data and structural images, respectively. Here $f_{G}$ can be any geometry-based model (\eg, SchNet~\cite{schutt2017schnet}, SphereNet~\cite{liu2022spherical}, Equiformer~\cite{liao2023equiformer}, ViSNet~\cite{wang2024enhancing}).

A mapper $f_M$ (a single linear layer) projects geometry features into the structural space; the nonlinear transformation is handled by the frozen ED predictor $f_{EDP}$:
\begin{align}
\label{Formula_ED_features}
\mathcal{F}^{\mathcal{G}\rightarrow\mathcal{U}}=f_{EDP}(f_M(\mathcal{F}^{\mathcal{G}})); \mathcal{F}^{\mathcal{S}\rightarrow\mathcal{U}}=f_{EDP}(\mathcal{F}^{\mathcal{S}})
\end{align}

\textbf{EDG: Naive Distillation.} In the base framework EDG, we distill knowledge from all samples:
\begin{equation}
    \mathcal{L}_{ED}^{naive} = SL1(\mathcal{F}^{\mathcal{G}\rightarrow\mathcal{U}}, \mathcal{F}^{\mathcal{S}\rightarrow\mathcal{U}})
\end{equation}
where $SL1$ represents smooth L1 distance~\cite{girshick2015fast}.

\textbf{EDG++: Selective Distillation.} Given the adaptive threshold $\tau$, the selection mask $m_i = \mathbf{1}[c_i \geq \tau]$ determines which samples receive distillation supervision. The selective distillation loss is:
\begin{equation}
    \mathcal{L}_{ED}^{selective} = \begin{cases} \frac{\sum_{i} m_i \cdot SL1(\mathcal{F}^{\mathcal{G}\rightarrow\mathcal{U}}_i, \mathcal{F}^{\mathcal{S}\rightarrow\mathcal{U}}_i)}{\sum_i m_i}, & \text{if } \sum_i m_i > 0 \\ 0, & \text{otherwise} \end{cases}
\end{equation}
When no samples pass the threshold, only the task loss drives optimization.

\textbf{Total Training Loss.} The task predictor $f_T$ outputs predictions $\hat{y}=f_T(\mathcal{F}^{\mathcal{G}})$, and the final loss is:
\begin{align}
\label{Formula_L_EDG}
\mathcal{L}_{EDG}=\underbrace{L1(\hat{y}, y)}_{\mathcal{L}_{Task}}+\lambda\mathcal{L}_{ED}
\end{align}
where $\lambda$ controls the distillation strength, and $\mathcal{L}_{ED}$ is either $\mathcal{L}_{ED}^{naive}$ (EDG) or $\mathcal{L}_{ED}^{selective}$ (EDG++).

\textbf{Inference.} At inference time, only the geometry student and task predictor are used. The ED-aware teacher, reliability estimator, and distillation-related components are not required, ensuring no computational overhead compared to baseline geometry models.
